\section{Task and Evaluation Model}
\label{sec:method}

We formalize DeFi workflow synthesis as mapping a natural-language intent
$p$ to a workflow specification $W = (N, E, C)$, where $N$ is a set of
typed nodes drawn from a fixed DeFi node vocabulary (e.g.\
\texttt{tokenSelector}, \texttt{oneInchQuote}, \texttt{priceImpact\-Calculator},
\texttt{oneInchSwap}, \texttt{limitOrder}, \texttt{fusionPlus},
\texttt{transactionMonitor}), $E \subseteq N \times N$ is a set of
type-level data-flow edges, and $C$ is an execution configuration mapping
keys such as \texttt{fromToken}, \texttt{amount}, \texttt{slippage}, and
\texttt{targetPrice} to values.

Each benchmark item pairs a prompt with an expert-authored gold
annotation $(N^{\star}, E^{\star}, C^{\star}, S^{\star})$, where
$S^{\star}$ is a set of \emph{safety predicates} the workflow must satisfy
to be safely executable (e.g.\ a slippage bound, a price-impact gate,
transaction monitoring, a distinct bridge destination). The annotation
also marks whether a prompt is underspecified, in which case the ideal
system response is to request clarification rather than emit a workflow.

\subsection{Three static evaluation layers}

We score a generated workflow along three strictly nested layers, so that
each captures a distinct failure mode.

\paragraph{Structural validity.} A workflow is \emph{graph-valid} if it
recalls every required node type and every required type-level edge and
the generator did not error. This layer is deliberately lenient: it
ignores extra nodes unless they break a required edge, matching the weak
notion of ``valid'' that graph-only evaluations reward.

\paragraph{Executable completeness.} A workflow is \emph{executable} if it
is graph-valid \emph{and} every required configuration key in $C^{\star}$
is populated with a concrete (non-placeholder) value. A node that merely
exists in template mode, with no trade amount or target price, is not
executable.

\paragraph{Declared safety.} A workflow is \emph{statically safe} if it is
executable \emph{and} every required safety predicate in $S^{\star}$ is
declared and parameterized. Safety predicates are derived uniformly from
the normalized workflow for all systems, so no system is advantaged by
self-reporting safety.

A separate \emph{clarification} axis scores underspecified prompts: the
correct behavior is to ask for the missing information, not to emit a
confidently wrong workflow.

\subsection{On-chain execution layer}
\label{sec:method-onchain}

Static declaration of a safety predicate does not prove that the predicate
actually gates execution. We therefore add an execution layer that runs
each generated swap on a real EVM. We deploy a constant-product
automated market maker (Uniswap-V2-style $x\cdot y = k$ with a $0.3\%$ fee)
and mock ERC20 tokens on an in-process EVM, seed deterministic reserves,
and submit the generated trade as an actual transaction. The harness
observes the mined receipt, gas, or revert, and labels the outcome as one
of \texttt{executed\_safe}, \texttt{reverted\_slippage},
\texttt{aborted\_price\_impact}, \texttt{unsafe\_executed},
\texttt{pending\_not\_filled}, or \texttt{not\_executable}.

The critical distinction the harness captures is between a slippage bound
and a price-impact gate. A slippage bound sets a minimum output computed
from the (already impact-adjusted) quote, so it protects only against
price \emph{movement between quote and execution}; it does \emph{not}
prevent a large order from executing at severe self-inflicted price
impact. Only a separate price-impact gate does. A workflow that mines a
high-impact trade with no such gate is labeled \texttt{unsafe\_executed}.
This is a deterministic local EVM snapshot, not a mainnet fork: it
provides real EVM execution semantics (real \texttt{transferFrom}, real
\texttt{require} reverts, real gas) but synthetic prices and liquidity
(Section~\ref{sec:limitations}).
